% !TEX program = xelatex
\documentclass[11pt,a4paper]{article}
\usepackage{algo-preamble}

\title{\textbf{L11 \textperiodcentered{} Άπληστοι Αλγόριθμοι I}\\[2pt]
\large Χρονοπρογραμματισμός Διαστημάτων}
\author{Algorithms Class Hub}
\date{\small Αλγόριθμοι και Πολυπλοκότητα (K17) \textperiodcentered{} ΕΚΠΑ}

\begin{document}
\maketitle

\begin{center}\itshape\small
Η ιδέα του άπληστου αλγορίθμου, ο χρονοπρογραμματισμός διαστημάτων (interval
scheduling) με το κριτήριο του χρόνου λήξης, η απόδειξη ορθότητας, και η
διαμέριση διαστημάτων.
\end{center}

\section{Τι θα δούμε}

Ξεκινάει μια νέα οικογένεια αλγορίθμων: οι \textbf{άπληστοι} (greedy). Ήδη
συναντήσαμε τρεις στο μάθημα για τα γραφήματα --- Dijkstra, Prim, Kruskal ---
αλλά εκεί η απληστία ήταν «κρυμμένη» μέσα στο γράφημα. Εδώ τη βλέπουμε καθαρά, σε
δύο προβλήματα χρονοπρογραμματισμού:
\begin{enumerate}
  \item \textbf{Χρονοπρογραμματισμός διαστημάτων} --- διάλεξε όσο το δυνατόν
  περισσότερες μη-επικαλυπτόμενες εργασίες.
  \item \textbf{Διαμέριση διαστημάτων} --- μοίρασε όλες τις εργασίες σε όσο το
  δυνατόν λιγότερες μηχανές.
\end{enumerate}

Το πιο σημαντικό μάθημα εδώ \textbf{δεν} είναι ο αλγόριθμος --- είναι το πώς
\textbf{αποδεικνύεις} ότι ένας άπληστος αλγόριθμος είναι σωστός. Θα το χρειαστείς
ξανά και ξανά στην εξεταστική.

\section{Η ιδέα του άπληστου αλγορίθμου}

\begin{keypoint}
\textbf{Άπληστος αλγόριθμος:} χτίζει τη λύση \textbf{σταδιακά}, και σε κάθε βήμα
παίρνει την απόφαση που φαίνεται \textbf{τοπικά καλύτερη} --- χωρίς να
ξανασκέφτεται, χωρίς να κοιτάει το μέλλον.
\end{keypoint}

\begin{intuition}
\textbf{Διαισθητικά:} ο άπληστος αλγόριθμος είναι σαν να ανεβαίνεις βουνό
κοιτάζοντας μόνο τα πόδια σου --- κάθε βήμα προς τα πάνω. Μερικές φορές αυτό σε
φέρνει στην κορυφή· άλλες φορές σε παγιδεύει σε έναν μικρό λοφίσκο. \textbf{Το
ζητούμενο σε κάθε greedy πρόβλημα είναι το ίδιο: ποιο τοπικό κριτήριο εγγυάται
την παγκόσμια βέλτιστη λύση --- και γιατί;}
\end{intuition}

Δεν αρκεί να βρεις έναν άπληστο κανόνα που «δείχνει λογικός». Πρέπει να τον
\textbf{αποδείξεις}. Όπως θα δούμε αμέσως, οι περισσότεροι λογικοφανείς κανόνες
είναι \textbf{λάθος}.

\section{Χρονοπρογραμματισμός Διαστημάτων (Interval Scheduling)}

\subsection{Διατύπωση}

\textbf{Είσοδος:} $n$ εργασίες. Η εργασία $j$ ξεκινά τη χρονική στιγμή $s_j$ και
τελειώνει τη στιγμή $f_j$.

\textbf{Συμβατές εργασίες:} δύο εργασίες είναι \textbf{συμβατές} αν δεν
επικαλύπτονται χρονικά.

\textbf{Στόχος:} βρες το \textbf{μεγαλύτερο σε πλήθος} υποσύνολο εργασιών που
είναι ανά δύο συμβατές.

\begin{intuition}
\textbf{Πραγματικό παράδειγμα:} μία αίθουσα, πολλές αιτήσεις για κρατήσεις,
καθεμία με ώρα έναρξης και λήξης. Θέλεις να ικανοποιήσεις \textbf{όσο το δυνατόν
περισσότερες} --- όχι τις «μεγαλύτερες» ή τις «σημαντικότερες», απλώς το μέγιστο
πλήθος.
\end{intuition}

\subsection{Το άπληστο πρότυπο}

Όλοι οι υποψήφιοι αλγόριθμοι έχουν την ίδια σκελετική μορφή:

\begin{quoteblock}
Θεώρησε τις εργασίες σε \textbf{κάποια σειρά}. Διάλεξε κάθε εργασία εφόσον είναι
συμβατή με όσες έχεις ήδη επιλέξει.
\end{quoteblock}

Όλο το παιχνίδι είναι \textbf{ποια σειρά}. Τέσσερις λογικοφανείς υποψήφιες:

\begin{center}
\begin{tabular}{@{}l l@{}}
\toprule
Κριτήριο & Σειρά \\
\midrule
Μικρότερος χρόνος έναρξης & αύξουσα ως προς $s_j$ \\
Μικρότερος χρόνος λήξης   & αύξουσα ως προς $f_j$ \\
Μικρότερο διάστημα        & αύξουσα ως προς $f_j - s_j$ \\
Ελάχιστες διενέξεις       & αύξουσα ως προς $c_j$ (πλήθος μη-συμβατών εργασιών) \\
\bottomrule
\end{tabular}
\end{center}

\subsection{Τρία κριτήρια που αποτυγχάνουν}

\begin{warningbox}
\textbf{Μικρότερος χρόνος έναρξης --- ΛΑΘΟΣ.} Η εργασία που ξεκινάει πρώτη
μπορεί να είναι τεράστια και να επικαλύπτει πολλές άλλες, μπλοκάροντας μια λύση
με περισσότερες εργασίες.

\textbf{Μικρότερο διάστημα --- ΛΑΘΟΣ.} Η συντομότερη εργασία μπορεί να «κάθεται»
ακριβώς ανάμεσα σε δύο εργασίες που μεταξύ τους είναι συμβατές --- διαλέγοντάς
την, χάνεις και τις δύο.

\textbf{Ελάχιστες διενέξεις --- ΛΑΘΟΣ.} Με ανάλογο αντιπαράδειγμα, η εργασία με
τις λιγότερες συγκρούσεις μπορεί και πάλι να μπλοκάρει μια μεγαλύτερη λύση.
\end{warningbox}

\textbf{Αντιπαράδειγμα για το «μικρότερο διάστημα».} Έστω τρεις εργασίες: η $A$
(διάστημα 5) και η $C$ (διάστημα 4) είναι μεταξύ τους συμβατές, ενώ η $B$
(διάστημα 2) κάθεται ανάμεσά τους και συγκρούεται και με τις δύο. Το κριτήριο
«μικρότερο διάστημα» διαλέγει πρώτο το $B$ --- και τότε ούτε το $A$ ούτε το $C$
είναι πια συμβατά. Λύση: $1$. Η βέλτιστη λύση $\{A, C\}$ έχει μέγεθος $2$.

\subsection{Το κριτήριο που δουλεύει: μικρότερος χρόνος λήξης}

\begin{keypoint}
\textbf{Άπληστος αλγόριθμος (earliest finish time).} Θεώρησε τις εργασίες σε
\textbf{αύξουσα σειρά χρόνου λήξης}. Διάλεξε κάθε εργασία εφόσον είναι συμβατή με
τις ήδη επιλεγμένες.
\end{keypoint}

\begin{intuition}
\textbf{Γιατί ο χρόνος λήξης;} Επειδή η εργασία που τελειώνει \textbf{νωρίτερα}
αφήνει \textbf{όσο το δυνατόν περισσότερο ελεύθερο χρόνο} για όλες τις επόμενες.
Είναι η πιο «γενναιόδωρη» επιλογή προς το μέλλον --- και αυτή ακριβώς η διαίσθηση
γίνεται η απόδειξη παρακάτω.
\end{intuition}

\begin{codebox}
\begin{verbatim}
ΧρονοπρογραμματισμόςΔιαστημάτων(s, f):
  Ταξινόμησε τις εργασίες ώστε f_1 <= f_2 <= ... <= f_n
  A <- ∅
  Για j = 1 έως n:
    Εάν η εργασία j είναι συμβατή με το A:
      A <- A ∪ {j}
  επίστρεψε A
\end{verbatim}
\end{codebox}

\textbf{Υλοποίηση σε $O(n \log n)$:} η ταξινόμηση κοστίζει $O(n \log n)$. Για τον
έλεγχο συμβατότητας δεν χρειάζεται να κοιτάμε όλο το $A$ --- αρκεί να θυμόμαστε
την \textbf{τελευταία} εργασία $j^*$ που προστέθηκε. Η εργασία $j$ είναι συμβατή
με το $A$ \textbf{αν και μόνο αν} $s_j \ge f_{j^*}$. Άρα κάθε έλεγχος είναι
$O(1)$.

(Σε ένα τυπικό στιγμιότυπο 8 εργασιών, ο άπληστος κατά χρόνο λήξης μπορεί να
επιλέξει π.χ.\ το σύνολο $\{1, 4, 8\}$ μεγέθους 3 --- που είναι το βέλτιστο.)

\subsection{Ορθότητα --- «ο άπληστος προηγείται»}

\begin{keypoint}
\textbf{Θεώρημα.} Ο άπληστος αλγόριθμος ως προς τον χρόνο λήξης επιστρέφει λύση
\textbf{μέγιστου} πλήθους.
\end{keypoint}

\textbf{Απόδειξη (εις άτοπον).} Έστω $i_1, i_2, \dots, i_k$ οι εργασίες που
επιλέγει ο άπληστος, σε σειρά. Έστω $j_1, j_2, \dots, j_m$ μια \textbf{βέλτιστη}
λύση, διαλεγμένη έτσι ώστε να \textbf{συμφωνεί με τον άπληστο για όσο το δυνατόν
περισσότερες αρχικές εργασίες} --- δηλαδή $i_1 = j_1, \dots, i_r = j_r$ για τη
μέγιστη δυνατή τιμή του $r$. Αν ήταν $k < m$, ο άπληστος δεν είναι βέλτιστος. Θα
φτάσουμε σε άτοπο.

\textbf{Ο άπληστος «προηγείται».} Η εργασία $i_{r+1}$ που διάλεξε ο άπληστος μετά
την $i_r$ έχει χρόνο λήξης $f(i_{r+1}) \le f(j_{r+1})$. Γιατί; Επειδή
$i_r = j_r$, η εργασία $j_{r+1}$ είναι συμβατή με την $i_r$ --- άρα ήταν
\textbf{υποψήφια} όταν ο άπληστος διάλεγε την $i_{r+1}$. Και ο άπληστος επιλέγει
πάντα την υποψήφια με τον \textbf{μικρότερο} χρόνο λήξης. Άρα
$f(i_{r+1}) \le f(j_{r+1})$.

\textbf{Η ανταλλαγή.} Φτιάχνουμε μια νέα λύση αντικαθιστώντας στη βέλτιστη την
$j_{r+1}$ με την $i_{r+1}$:
\[
S' = j_1, \dots, j_r,\ i_{r+1},\ j_{r+2}, \dots, j_m.
\]
Είναι έγκυρη: η $i_{r+1}$ είναι συμβατή με τις $j_1, \dots, j_r$
($= i_1, \dots, i_r$) αφού ο άπληστος την επέλεξε, και επειδή
$f(i_{r+1}) \le f(j_{r+1})$ τελειώνει \textbf{όχι αργότερα} από την $j_{r+1}$,
άρα δεν συγκρούεται ούτε με την $j_{r+2}$. Το $S'$ έχει $m$ εργασίες --- άρα
είναι \textbf{κι αυτό βέλτιστο} --- αλλά συμφωνεί με τον άπληστο για $r+1$
εργασίες. \textbf{Αντίφαση} στη μεγιστότητα του $r$.

Μένει η περίπτωση $k = r < m$: τότε υπάρχει η $j_{r+1}$, συμβατή με την $i_r$,
και ο άπληστος \textbf{δεν} θα είχε σταματήσει --- θα την είχε εξετάσει και θα
είχε προσθέσει κάτι. Άτοπο. Άρα $k = m$. $\blacksquare$

\begin{intuition}
\textbf{Το όνομα της τεχνικής:} «ο άπληστος προηγείται» (greedy stays ahead).
Δείχνεις ότι μετά από κάθε βήμα, η μερική λύση του άπληστου είναι
\textbf{τουλάχιστον τόσο καλή} όσο οποιασδήποτε άλλης λύσης --- εδώ, ότι
τελειώνει νωρίτερα. Θα δούμε άλλες δύο τεχνικές απόδειξης στο επόμενο μάθημα.
\end{intuition}

\section{Διαμέριση Διαστημάτων (Interval Partitioning)}

\subsection{Διατύπωση}

Αλλάζουμε το ερώτημα. Τώρα \textbf{πρέπει} να προγραμματίσουμε \textbf{όλες} τις
εργασίες.

\textbf{Είσοδος:} $n$ εργασίες, η $j$ με $(s_j, f_j)$.

\textbf{Στόχος:} βρες το \textbf{ελάχιστο πλήθος μηχανών} ώστε να αναθέσεις κάθε
εργασία σε μία μηχανή, με δύο εργασίες που τρέχουν ταυτόχρονα \textbf{ποτέ} στην
ίδια μηχανή.

\begin{intuition}
\textbf{Πραγματικό παράδειγμα:} μαθήματα που πρέπει όλα να γίνουν, και ψάχνεις
τον \textbf{ελάχιστο αριθμό αιθουσών}. Δύο μαθήματα που πέφτουν την ίδια ώρα
χρειάζονται διαφορετική αίθουσα.
\end{intuition}

\subsection{Το κάτω φράγμα: το βάθος}

\begin{keypoint}
\textbf{Βάθος} ενός συνόλου διαστημάτων: το \textbf{μέγιστο πλήθος} διαστημάτων
που περιέχουν μια \textbf{κοινή} χρονική στιγμή.
\end{keypoint}

Αν σε κάποια στιγμή τρέχουν ταυτόχρονα $d$ εργασίες, αυτές χρειάζονται $d$
διαφορετικές μηχανές. Άρα:
\[
\text{πλήθος μηχανών} \ \ge\ \text{βάθος}.
\]
Το ερώτημα είναι: \textbf{μπορούμε πάντα να πετύχουμε ακριβώς το βάθος;}

\subsection{Ο άπληστος αλγόριθμος}

\begin{keypoint}
\textbf{Άπληστο κριτήριο.} Θεώρησε τις εργασίες σε \textbf{αύξουσα σειρά χρόνου
έναρξης}. Ανάθεσε κάθε εργασία σε μια \textbf{οποιαδήποτε ελεύθερη} μηχανή· αν
δεν υπάρχει, \textbf{άνοιξε} μια νέα.
\end{keypoint}

\begin{codebox}
\begin{verbatim}
ΔιαμέρισηΔιαστημάτων(s, f):
  Ταξινόμησε τις εργασίες ώστε s_1 <= s_2 <= ... <= s_n
  d <- 0                                  // πλήθος μηχανών
  Για j = 1 έως n:
    Εάν η εργασία j είναι συμβατή με κάποια ήδη ανοιγμένη μηχανή k:
      ανάθεσε την j στη μηχανή k
    Αλλιώς:
      d <- d + 1
      άνοιξε τη μηχανή d και ανάθεσε την j σε αυτήν
  επίστρεψε d
\end{verbatim}
\end{codebox}

\subsection{Ορθότητα}

\textbf{Παρατήρηση πρώτα:} ο αλγόριθμος δεν αναθέτει \textbf{ποτέ} δύο
συγκρουόμενες εργασίες στην ίδια μηχανή --- εξ ορισμού αναθέτει σε μηχανή μόνο αν
είναι «συμβατή» με ό,τι έχει ήδη εκεί. Άρα η λύση είναι \textbf{έγκυρη}. Μένει να
δείξουμε ότι χρησιμοποιεί \textbf{ελάχιστες} μηχανές.

\begin{keypoint}
\textbf{Θεώρημα.} Ο άπληστος αλγόριθμος είναι βέλτιστος: χρησιμοποιεί ακριβώς
$d = \text{βάθος}$ μηχανές.
\end{keypoint}

\textbf{Απόδειξη.} Έστω $d$ το πλήθος των μηχανών που άνοιξε ο αλγόριθμος. Η
τελευταία μηχανή, η $d$-οστή, άνοιξε επειδή κάποια εργασία --- έστω η $j$ ---
\textbf{δεν} ήταν συμβατή με καμία από τις $d-1$ μηχανές που ήδη υπήρχαν.

Άρα, τη στιγμή που εξετάστηκε η $j$, καθεμία από τις $d-1$ μηχανές είχε
\textbf{τουλάχιστον μία} εργασία σε σύγκρουση με την $j$. Επειδή ο αλγόριθμος
εξετάζει τις εργασίες σε \textbf{αύξουσα σειρά έναρξης}, όλες αυτές οι εργασίες
ξεκίνησαν \textbf{πριν} την $s_j$ --- και αφού συγκρούονται με την $j$, δεν έχουν
τελειώσει μέχρι την $s_j$.

Συνεπώς, τη χρονική στιγμή $s_j + \epsilon$ τρέχουν ταυτόχρονα \textbf{$d$
εργασίες} (οι $d-1$ συν την $j$). Άρα το βάθος είναι $\ge d$. Αλλά κάθε λύση
χρειάζεται $\ge \text{βάθος} \ge d$ μηχανές. Ο άπληστος χρησιμοποιεί ακριβώς $d$.
Άρα είναι βέλτιστος. $\blacksquare$

\begin{intuition}
\textbf{Το όνομα της τεχνικής:} \textbf{δομική εξήγηση} (structural bound).
Βρήκαμε ένα «εμπόδιο» μέσα στο ίδιο το πρόβλημα --- το βάθος --- που κανένας
αλγόριθμος δεν μπορεί να ξεπεράσει, και δείξαμε ότι ο άπληστος ακριβώς αυτό το
εμπόδιο πετυχαίνει. Όταν το κάτω φράγμα συναντά την τιμή του αλγορίθμου,
ησυχάζεις.
\end{intuition}

\begin{recapbox}
\textbf{Τι κρατάμε από αυτή τη διάλεξη}
\begin{enumerate}
  \item \textbf{Άπληστος αλγόριθμος} --- χτίζει λύση σταδιακά, κάθε βήμα τοπικά
  βέλτιστο. Το δύσκολο δεν είναι ο κανόνας, είναι η \textbf{απόδειξη}.
  \item \textbf{Χρονοπρογραμματισμός διαστημάτων} --- μέγιστο πλήθος συμβατών
  εργασιών. Σωστό κριτήριο: \textbf{μικρότερος χρόνος λήξης}. $O(n \log n)$.
  \item Τα κριτήρια «μικρότερη έναρξη», «μικρότερο διάστημα», «ελάχιστες
  διενέξεις» \textbf{αποτυγχάνουν} --- πάντα ζήτα αντιπαράδειγμα πριν πιστέψεις
  έναν greedy κανόνα.
  \item \textbf{Greedy stays ahead} --- απόδειξη ότι η μερική λύση του άπληστου
  προηγείται κάθε άλλης σε κάθε βήμα.
  \item \textbf{Διαμέριση διαστημάτων} --- ελάχιστο πλήθος μηχανών. Κριτήριο:
  \textbf{μικρότερος χρόνος έναρξης}. Το αποτέλεσμα ισούται πάντα με το
  \textbf{βάθος}.
  \item \textbf{Structural bound} --- βρες ένα κάτω φράγμα μέσα στο πρόβλημα (το
  βάθος) και δείξε ότι ο αλγόριθμος το πετυχαίνει.
\end{enumerate}
\end{recapbox}

\lecturefooter
\end{document}
